\chapter[Metropolis--Hastings dhe Monte Carlo në fizikën statistike]{Metropolis--Hastings dhe Monte Carlo\\në fizikën statistike}
\begin{qellime}
Studenti ndërton algoritmin Metropolis--Hastings; provon ruajtjen e shpërndarjes synuar; dallon termalizimin, përzierjen dhe autokorrelacionin; simulon modelin Ising dhe interpreton energjinë, magnetizimin, kapacitetin termik dhe tranzicionin fazor.
\end{qellime}

\section{Pse nevojitet MCMC}
Në fizikën statistike duam pritje sipas shpërndarjes Boltzmann,
\begin{equation}
 \pi(x)=\frac{e^{-\beta E(x)}}{Z},\qquad
 Z=\sum_xe^{-\beta E(x)}.
\end{equation}
Numri i gjendjeve rritet eksponencialisht dhe $Z$ është zakonisht i paarritshëm. Metodat Markov Chain Monte Carlo ndërtojnë një zinxhir me shpërndarje stacionare $\pi$ pa kërkuar konstanten e normalizimit.

\section{Algoritmi Metropolis--Hastings}
Nga gjendja $x$, propozohet $y\sim q(y|x)$ dhe pranohet me probabilitet
\begin{equation}
 \alpha(x,y)=\min\left[1,\frac{\pi(y)q(x|y)}{\pi(x)q(y|x)}\right].
\end{equation}
Nëse propozimi është simetrik,
\begin{equation}
 \alpha(x,y)=\min[1,e^{-\beta\Delta E}].
\end{equation}
Raporti anulon $Z$. Hapi i refuzuar nuk hidhet; gjendja e vjetër regjistrohet sërish. Kjo është e domosdoshme për shpërndarjen e saktë.

\begin{shembull}[Algoritmi Metropolis--Hastings]
Niset nga $x_0$. Për çdo hap propozohet $y\sim q(\cdot|x_n)$ dhe gjenerohet $u\sim\mathcal U(0,1)$. Llogaritet $a=\min\{1,\pi(y)q(x_n|y)/[\pi(x_n)q(y|x_n)]\}$. Nëse $u<a$, vendoset $x_{n+1}=y$; përndryshe vendoset $x_{n+1}=x_n$.
\end{shembull}

\section{Termalizimi dhe diagnostika}
Hapat e parë varen nga gjendja fillestare dhe hiqen si \emph{burn-in}. Gjatësia nuk zgjidhet me një numër universal; kontrollohen disa nisje, gjurmët, observablat dhe koha e autokorrelacionit. Hollimi i zinxhirit zakonisht hedh informacion dhe nuk zëvendëson vlerësimin e $N_{\mathrm{eff}}$.

Norma e pranimit nuk është qëllim në vetvete. Hapa shumë të vegjël pranohen shpesh, por lëvizin ngadalë; hapa shumë të mëdhenj refuzohen. Propozimi duhet të balancojë lëvizjen dhe pranimin. Për shpërndarje shumëmodale, zinxhiri mund të mbetet në një modë edhe kur gjurma duket e qëndrueshme.

\section{Modeli Ising dydimensional}
Në rrjet katror me spine $s_i=\pm1$,
\begin{equation}
 E=-J\sum_{\langle i,j\rangle}s_is_j-H\sum_is_i,
\qquad m=\frac1N\sum_is_i.
\end{equation}
Me kufij periodikë, çdo spin ka katër fqinjë. Për një kthim $s_i\to-s_i$,
\begin{equation}
 \Delta E=2s_i\left(J\sum_{j\in\mathrm{fqinj}}s_j+H\right).
\end{equation}
Në temperaturë të ulët, gjendjet e rreshtuara dominojnë; në temperaturë të lartë, entropia prodhon çrregullim.

\begin{figure}[ht]
\centering\includegraphics[width=.78\textwidth]{24_ising_konfigurime.pdf}
\caption{Formimi i domeneve nën temperaturën kritike dhe çrregullimi sipër saj.}
\end{figure}

\section{Observablat termodinamike}
Nga luhatjet merren
\begin{equation}
 C_V=\frac{\beta^2}{N}(\avg{E^2}-\avg E^2),
 \qquad
 \chi=\frac{\beta}{N}(\avg{M^2}-\avg M^2).
\end{equation}
Në sistem të fundmë, tranzicioni shfaqet si maksimum i gjerë, jo singularitet i vërtetë. Temperatura kritike e rrjetit katror pa fushë është
\begin{equation}
 \frac{k_BT_c}{J}=\frac{2}{\ln(1+\sqrt2)}\approx2.269.
\end{equation}

\begin{figure}[ht]
\centering\includegraphics[width=.74\textwidth]{25_ising_magnetizimi.pdf}
\caption{Rënia e magnetizimit pranë zonës kritike në një simulim të fundmë.}
\end{figure}

\subsection{Ngadalësimi kritik dhe përmirësimet}
Pranë $T_c$, domenet ekzistojnë në shumë shkallë dhe përditësimi me një spin përzihet ngadalë. Koha e autokorrelacionit rritet me madhësinë e sistemit. Algoritmet e grupeve, si Wolff dhe Swendsen--Wang, ndryshojnë bashkë rajone të korreluara dhe ulin ngadalësimin kritik.

\subsection{Gabimet dhe analiza e të dhënave MCMC}
Mesatarja llogaritet pas termalizimit, por gabimi nuk merret duke trajtuar çdo hap si të pavarur. Përdoren autokorrelacioni i integruar, bllokimi ose batch means. Simulime të pavarura me fara të ndryshme ndihmojnë të zbulohet mos-përzierja.

\begin{kujdes}
Një konfigurim tërheqës vizualisht nuk është rezultat termodinamik. Përfundimet kërkojnë seri të termalizuara, madhësi efektive, prova me përmasa të ndryshme të rrjetit dhe observabla me gabime statistike.
\end{kujdes}

\subsection{Nga Ising te probleme të tjera}
I njëjti parim përdoret në mekanikën statistike, inferencën Bayesiane, modelet e materialeve dhe fushat në rrjet. Ndryshojnë energjia, hapësira e gjendjeve dhe propozimi; mbeten invarianca e shpërndarjes, ergodiciteti dhe analiza e korrelacionit.

\begin{lidhje}
Laboratori zbaton Ising 2D me Metropolis, verifikon $\Delta E$, mat normën e pranimit, termalizimin dhe autokorrelacionin, pastaj paraqet $|m|$, $C_V$ dhe $\chi$ kundrejt temperaturës.
\end{lidhje}

\section*{Pyetje kontrolli}
\begin{enumerate}
\item Pse konstanta e normalizimit anulohet në raportin Metropolis?
\item Pse hapat e refuzuar duhet të regjistrohen?
\item Çfarë është ngadalësimi kritik dhe si ndikon në gabimin statistik?
\item Pse një simulim me një madhësi rrjeti nuk mjafton për të vërtetuar tranzicion fazor?
\end{enumerate}
\section{Algoritmi Metropolis për modelin Ising}
\subsection{Energjia lokale dhe një sweep}
\begin{lstlisting}[language=Python,caption={Përditësimi Metropolis me një spin.}]
def delta_E(spin, i, j, J=1.0, H=0.0):
    L = spin.shape[0]
    fq = (spin[(i+1) % L, j] + spin[(i-1) % L, j]
          + spin[i, (j+1) % L] + spin[i, (j-1) % L])
    return 2.0*spin[i, j]*(J*fq + H)

def sweep_metropolis(spin, beta, rng, J=1.0, H=0.0):
    L = spin.shape[0]
    pranuar = 0
    for _ in range(L*L):
        i, j = rng.integers(0, L, size=2)
        dE = delta_E(spin, i, j, J, H)
        if dE <= 0.0 or rng.random() < np.exp(-beta*dE):
            spin[i, j] *= -1
            pranuar += 1
    return pranuar/(L*L)
\end{lstlisting}

\subsection{Matja e energjisë pa numërim të dyfishtë}
\begin{lstlisting}[language=Python]
def observablat(spin, J=1.0, H=0.0):
    djathtas = np.roll(spin, -1, axis=1)
    poshte = np.roll(spin, -1, axis=0)
    E = -J*np.sum(spin*(djathtas + poshte)) - H*np.sum(spin)
    M = np.sum(spin)
    return E, M
\end{lstlisting}
Numërohen vetëm lidhjet djathtas dhe poshtë, në mënyrë që çdo çift fqinjësh të hyjë një herë.

\subsection{Termalizimi dhe matjet}
\begin{lstlisting}[language=Python,caption={Struktura e një simulimi të plotë.}]
def simulo_ising(L, T, n_term, n_matje, hendek, seed):
    rng = np.random.default_rng(seed)
    spin = rng.choice([-1, 1], size=(L, L))
    beta = 1.0/T
    for _ in range(n_term):
        sweep_metropolis(spin, beta, rng)
    E, M = [], []
    for _ in range(n_matje):
        for _ in range(hendek):
            sweep_metropolis(spin, beta, rng)
        e, m = observablat(spin)
        E.append(e); M.append(m)
    return np.asarray(E), np.asarray(M)
\end{lstlisting}
Hendeku nuk garanton pavarësi. Autokorrelacioni matet dhe raportohet $N_{\mathrm{eff}}$.

\subsection{Observablat dhe gabimet}
\begin{lstlisting}[language=Python]
N = L*L
e = E/N; m = M/N
Cv = (np.mean(E**2) - np.mean(E)**2)/(N*T*T)
chi = (np.mean(M**2) - np.mean(M)**2)/(N*T)
m_abs = np.mean(np.abs(m))
\end{lstlisting}
Për gabimet përdoret bllokimi: seria ndahet në blloqe më të gjata se koha e korrelacionit dhe mesataret e blloqeve trajtohen si afërsisht të pavarura.

\subsection{Protokolli i studimit termik}
\begin{enumerate}
\item Ekzekutohen disa temperatura larg dhe pranë $T_c$.
\item Për çdo temperaturë kontrollohen termalizimi dhe pranimi.
\item Raportohen $\avg{|m|}$, $\avg e$, $C_V$, $\chi$ dhe gabimet.
\item Përsëritet për disa $L$ për të dalluar efektet e përmasës së fundme.
\item Krahasohen nisja e rregullt dhe e çrregullt pranë zonës kritike.
\end{enumerate}

\begin{kujdes}
Përdoren \emph{balancë e detajuar}, \emph{shpërndarje stacionare}, \emph{termalizim}, \emph{përzierje}, \emph{kohë autokorrelacioni} dhe \emph{ngadalësim kritik}. Këto terma përkufizohen para përdorimit dhe nuk zëvendësohen me anglicizma të panevojshëm.
\end{kujdes}

Termat e fizikës statistike janë harmonizuar me përdorimin universitar në shqip dhe me burimet probabilitare \cite{capollari,butka}; algoritmi dhe analiza MCMC mbështeten te burimet themelore \cite{metropolis,hastings,krauth,robert,walter}.

\subsection{Ndërtimi i Metropolis--Hastings nga balanca e detajuar}
Synohet shpërndarja $\pi(x)$, e njohur deri në një konstante normalizimi. Nga gjendja $x$ propozohet $y\sim q(y|x)$ dhe pranohet me probabilitet
\begin{equation}
 \alpha(x,y)=\min\left\{1,
 \frac{\pi(y)q(x|y)}{\pi(x)q(y|x)}\right\}.
\end{equation}
Për $x\ne y$, kerneli i kalimit është $T(x,y)=q(y|x)\alpha(x,y)$. Nëse raporti është më i vogël se një,
\begin{equation}
 \pi(x)T(x,y)=\pi(x)q(y|x)
 \frac{\pi(y)q(x|y)}{\pi(x)q(y|x)}
 =\pi(y)q(x|y),
\end{equation}
ndërsa kalimi i kundërt pranohet me probabilitet një. Rasti tjetër është simetrik; kështu vlen balanca e detajuar. Probabiliteti i refuzimit vendos masën e mbetur në $T(x,x)$ dhe siguron normalizimin.

Kur $q(y|x)=q(x|y)$, raporti i propozimit anulohet. Për shpërndarjen Boltzmann $\pi(x)\propto e^{-\beta E(x)}$ merret
\begin{equation}
 \alpha=\min\{1,e^{-\beta\Delta E}\}.
\end{equation}
Konstantja e ndarjes nuk llogaritet: ky është një nga avantazhet kryesore të MCMC.

\subsection{Zgjedhja e propozimit, pranimi dhe përzierja}
Një hap shumë i vogël pranohet shpesh, por lëviz ngadalë; një hap shumë i madh refuzohet shpesh. Norma e pranimit nuk është objektivi i vetëm. Duhet parë largësia e përshkuar, autokorrelacioni dhe $N_{\mathrm{eff}}$ për koston llogaritëse. Propozimi duhet të ruajë mbështetjen e synuar dhe ta bëjë zinxhirin të pakthyeshëm e aperiodik.

Termalizimi hiqet vetëm pas kontrollit të trajektoreve nga gjendje të shpërndara. Pastaj maten observablat në pjesën stacionare. Përfundimi shoqërohet me kohën e integruar të autokorrelacionit, madhësinë efektive dhe, kur është e mundur, krahasimin e shumë zinxhirëve.

\subsection{Modeli Ising dydimensional}
Në një rrjet katror me $N=L^2$ spine $s_i=\pm1$ dhe kufij periodikë,
\begin{equation}
 E=-J\sum_{\langle i,j\rangle}s_is_j-H\sum_i s_i,
 \qquad M=\sum_i s_i.
\end{equation}
Çdo lidhje numërohet një herë. Kur përmbyset spini $s_i$, vetëm lidhjet lokale ndryshojnë dhe
\begin{equation}
 \Delta E=2s_i\left(J\sum_{j\in\mathrm{fq}(i)}s_j+H\right).
\end{equation}
Ky rezultat shmang rillogaritjen e energjisë së gjithë rrjetit dhe ul koston e një prove nga $\Oh(N)$ në $\Oh(1)$. Një sweep përmban $N$ prova, pra kushton $\Oh(N)$.

Observablat termodinamike për spin janë
\begin{equation}
 e=\frac{\langle E\rangle}{N},\qquad
 m=\frac{\langle |M|\rangle}{N},
\end{equation}
\begin{equation}
 C_V=\frac{\beta^2}{N}\left(\langle E^2\rangle-\langle E\rangle^2\right),
 \qquad
 \chi=\frac{\beta}{N}\left(\langle M^2\rangle-\langle M\rangle^2\right).
\end{equation}
Përdorimi i $|M|$ në magnetizim shmang anulimin ndërmjet dy fazave të simetrisë në një rrjet të fundëm, por për susceptibilitetin duhet deklaruar qartë nëse përdoret $M$ apo $|M|$.

\subsection{Gabimet e matjeve të korreluara}
Mesatarja e energjisë nuk ka gabim $s_E/\sqrt{N_{\mathrm{matjeve}}}$ kur konfigurimet vijnë nga i njëjti zinxhir. Përdoret
\begin{equation}
 \operatorname{SE}(\overline E)=
 \sqrt{\frac{2\tau_{\mathrm{int},E}\Var(E)}{N_{\mathrm{matjeve}}}}.
\end{equation}
Për $C_V$ dhe $\chi$, që janë funksione jolineare të momenteve, bllokimi ose jackknife ruan korrelacionet dhe përhap pacaktueshmërinë. Madhësia e bllokut duhet të jetë dukshëm më e madhe se koha e autokorrelacionit; qëndrueshmëria e gabimit ndaj rritjes së bllokut është diagnostikë.

\subsection{Ngadalësimi kritik dhe madhësia e fundme}
Pranë temperaturës kritike, gjatësia e korrelacionit rritet dhe përditësimet lokale ndryshojnë vetëm ngadalë domenet e mëdha. Koha e autokorrelacionit rritet afërsisht si $\tau\sim L^z$. Ky ngadalësim kritik ul rëndë $N_{\mathrm{eff}}$. Algoritmet e grupimeve Wolff ose Swendsen--Wang përmbysin rajone të korreluara dhe mund ta zvogëlojnë këtë efekt, por Metropolis-i mbetet themeli pedagogjik për balancën e detajuar.

Për rrjet të fundëm nuk ka singularitet të vërtetë. Kulmet e $C_V$ dhe $\chi$ zhvendosen dhe mprehen me $L$. Prandaj një temperaturë kritike e vlerësuar nga një rrjet i vetëm është rezultat me efekt madhësie të fundme, jo kufiri termodinamik.

\begin{lstlisting}[language=Python,caption={Një sweep Metropolis për modelin Ising.}]
import numpy as np

def sweep_metropolis(spin, beta, rng, J=1.0, H=0.0):
    L = spin.shape[0]
    for _ in range(L * L):
        i = rng.integers(L); j = rng.integers(L)
        fqinjet = (spin[(i+1) % L, j] + spin[(i-1) % L, j]
                 + spin[i, (j+1) % L] + spin[i, (j-1) % L])
        delta_E = 2.0 * spin[i, j] * (J * fqinjet + H)
        if delta_E <= 0.0 or rng.random() < np.exp(-beta * delta_E):
            spin[i, j] *= -1
    return spin
\end{lstlisting}

\subsection{Protokolli i një studimi të riprodhueshëm}
Një studim termik duhet të deklarojë $L$, $J$, $H$, temperaturat, farat, gjendjet fillestare, numrin e sweeps të termalizimit, numrin e matjeve dhe mënyrën e vlerësimit të autokorrelacionit. Duhet provuar energjia lokale ndaj rillogaritjes së plotë në rrjete të vogla, të krahasohen kufijtë $T\to0$ dhe $T\to\infty$, dhe të kontrollohet simetria $M\leftrightarrow-M$ kur $H=0$.

\subsection{Përmbledhje dhe ushtrime}
Metropolis--Hastings ndërton një dinamikë artificiale me shpërndarje stacionare të zgjedhur. Koha Monte Carlo nuk është domosdoshmërisht kohë fizike; interpretimi termodinamik vjen nga shpërndarja e ekuilibrit.
\begin{enumerate}
\item Provoni balancën e detajuar për raportin e përgjithshëm Hastings.
\item Nxirrni $\Delta E$ për përmbysjen e një spini dhe kontrolloni numërimin e dyfishtë të lidhjeve.
\item Matni $\tau_{\mathrm{int}}$ të energjisë dhe magnetizimit në disa temperatura.
\item Krahasoni nisjen e renditur dhe të çrregullt për të përcaktuar termalizimin.
\item Kryeni studim me disa $L$ dhe diskutoni zhvendosjen e kulmit të susceptibilitetit.
\item Implementoni bllokimin ose jackknife për pacaktueshmërinë e kapacitetit termik.
\end{enumerate}

\subsection{Kufijtë analitikë të modelit Ising}
Në temperaturë zero dhe $H=0$, energjia minimizohet kur të gjitha spinet janë të njëjta. Në rrjetin katror çdo spin ka katër fqinjë, por çdo lidhje numërohet një herë; ka $2N$ lidhje. Prandaj energjia bazë është
\begin{equation}
 E_0=-2JN,qquad e_0=-2J.
\end{equation}
Magnetizimi për spin ka modul një. Këto janë teste të drejtpërdrejta për kodin në temperaturë të ulët.

Në temperaturë të pafundme, konfigurimet janë të barasvlershme dhe spinet janë të pavarura me mesatare zero. Atëherë $\langle E\rangle\to0$, $\langle M\rangle=0$ dhe
\begin{equation}
 \langle M^2\rangle=N.
\end{equation}
Prandaj $\langle M^2\rangle/N^2=1/N$. Një simulim me temperaturë shumë të lartë duhet t'i afrohet këtyre kufijve. Për rrjete shumë të vogla, të gjitha $2^N$ konfigurimet mund të enumerohen dhe të japin vlera të sakta të funksionit të ndarjes; ky është testi më i fortë i algoritmit dhe i observablave.

\subsubsection{Pranimi dhe bilanci energjetik lokal}
Në $H=0$, shuma e katër fqinjëve merr vlerat $-4,-2,0,2,4$. Për rrjedhojë, ndryshimi i energjisë kufizohet në bashkësinë
\begin{equation}
 \Delta E\in\{-8J,-4J,0,4J,8J\}.
\end{equation}
Faktorët $e^{-4\beta J}$ dhe $e^{-8\beta J}$ mund të parallogariten. Ky optimizim shmang eksponencialen në çdo provë dhe ruan saktësisht algoritmin.

Një matje e energjisë duhet të numërojë vetëm lidhjet djathtas dhe lart, ose të ndajë me dy shumën mbi të katër fqinjët. Gabimi i numërimit të dyfishtë ndryshon energjinë dhe temperaturën efektive. Testi lokal zgjedh një spin, llogarit energjinë e plotë para dhe pas përmbysjes dhe kontrollon barazinë me formulën e $\Delta E$ për shumë konfigurime të rastit.

\subsubsection{Kapaciteti termik, susceptibiliteti dhe cumulanti Binder}
Derivimi i pritjes kanonike jep
\begin{equation}
 \frac{\partial\langle E\rangle}{\partial\beta}
 =-\left(\langle E^2\rangle-\langle E\rangle^2\right).
\end{equation}
Me $\beta=1/(k_BT)$ merret formula e fluktuacioneve për $C_V$. Në mënyrë analoge, derivimi ndaj fushës jep susceptibilitetin nga fluktuacionet e magnetizimit. Këto marrëdhënie lidhin përgjigjen makroskopike me luhatjet mikroskopike.

Cumulanti Binder për magnetizimin është
\begin{equation}
 U_L=1-\frac{\langle M^4\rangle}{3\langle M^2\rangle^2}.
\end{equation}
Në fazën e çrregullt Gauss, $U_L\to0$; në një shpërndarje me dy maja të ngushta pranë $\pm M_0$, $U_L\to2/3$. Kurbat për madhësi të ndryshme rrjeti priten pranë temperaturës kritike, duke dhënë një vlerësim më robust se maksimumi i një observable të vetëm.

\subsubsection{Shkallëzimi me madhësinë e fundme}
Pranë pikës kritike, gjatësia e korrelacionit sillet si $\xi\sim|t|^{-\nu}$, ku $t=(T-T_c)/T_c$. Në rrjet të fundëm, rritja ndalet kur $\xi\sim L$. Kjo çon në forma shkallëzimi, për shembull
\begin{equation}
 m(T,L)=L^{-\beta_m/\nu}\mathcal M(tL^{1/\nu}),
\end{equation}
ku $\beta_m$ është eksponent kritik, jo inversi i temperaturës. Në $T_c$, grafiku log--log i $m$ ndaj $L$ ka pjerrësi $-\beta_m/\nu$.

Maksimumi i susceptibilitetit rritet afërsisht si $L^{\gamma/\nu}$ dhe temperatura e pseudokritike zhvendoset si $L^{-1/\nu}$. Një studim serioz përdor disa madhësi, interpolon observablat me pacaktueshmëri dhe kontrollon shembjen e të dhënave. Një kurbë e vetme magnetizimi nuk mjafton për të përcaktuar tranzicionin.

\subsubsection{Jackknife për madhësi jolineare}
Le të ndahen matjet në $B$ blloqe. Për çdo $b$ llogaritet statistiku $\theta_{(b)}$ duke hequr bllokun $b$. Mesatarja jackknife është $\overline\theta_J=B^{-1}\sum_b\theta_{(b)}$ dhe varianca vlerësohet me
\begin{equation}
 \widehat{\Var}_J(\theta)=\frac{B-1}{B}
 \sum_{b=1}^{B}(\theta_{(b)}-\overline\theta_J)^2.
\end{equation}
Kjo metodë përhap kovariancën ndërmjet $\langle E\rangle$ dhe $\langle E^2\rangle$ në $C_V$, ose ndërmjet momenteve në cumulantin Binder. Blloqet duhet të jenë më të gjata se autokorrelacioni; përndryshe pseudo-mostrat nuk janë afërsisht të pavarura.

\subsubsection{Krahasimi i algoritmeve}
Kostoja nuk matet vetëm me sweeps. Krahasimi përdor kohën e procesorit për një mostër efektivisht të pavarur,
\begin{equation}
 C_{\mathrm{eff}}\propto C_{\mathrm{sweep}}\,2\tau_{\mathrm{int}}.
\end{equation}
Metropolis-i lokal ka sweep të lirë, por $\tau$ të madh pranë kritikës. Algoritmi Wolff ndërton një grup spinesh të lidhura dhe e përmbys; kostoja për hap ndryshon, por autokorrelacioni kritik zvogëlohet fort. Krahasimi duhet bërë për të njëjtin observable dhe të njëjtën madhësi rrjeti.

Koha Monte Carlo e këtyre algoritmeve nuk është dinamika reale e magnetit, përveç rasteve kur rregullat e përditësimit janë zgjedhur si model kinetik. Për observablat e ekuilibrit mjafton shpërndarja stacionare; për koha relaksimi fizik nevojitet model dinamik me interpretim të veçantë.

